\section{Simulation Study}\label{sec:simulation}

The empirical analysis of this paper compares winter and summer spatial extremal dependence using pairwise nonparametric estimators. Before applying these estimators to the KNMI daily maximum wind-gust data, it is useful to study their finite-sample behaviour in a controlled setting where the true dependence structure is known. The purpose of the simulation study is therefore methodological: it assesses whether the proposed estimators are able to recover differences in spatial extremal dependence of the type expected in the empirical application. In particular, the simulation investigates whether the estimators can detect a regime in which winter extremes are more strongly and more widely spatially dependent than summer extremes.

The simulation is not intended to show that Dutch wind-gust extremes are generated by a specific parametric max-stable model. Instead, a max-stable benchmark is used as a controlled data-generating process. This allows the true pairwise extremal coefficient to be computed for every station pair and compared with the estimated dependence measures. The simulation therefore acts as a validation step for the empirical methodology: if the estimators fail to recover known dependence differences in finite samples comparable to the KNMI data, their interpretation in the empirical section would be questionable.

\subsection{Station geometry and sample size}\label{sec:simulation-geometry}

The simulated spatial fields are generated on the same type of station geometry as the empirical analysis. Let $(s_1,\ldots,s_d)$ denote the retained KNMI station locations, where $d$ is the number of stations in the fixed panel. For each station pair $(i,j)$, the inter-station distance is denoted by
\[
      h_{ij} = |s_i-s_j|,
\]
measured in kilometres. The use of the empirical station geometry ensures that the simulation reflects the same range of inter-station distances as the real-data analysis.

The simulated sample size is chosen to mimic the daily-maxima design. For each regime, a sample of $n$ spatial extreme fields is generated, where $n$ corresponds approximately to the number of seasonal daily maxima available in the empirical sample, pooled across all years. With 35 years of data and roughly 90 days per season, each station-season contains approximately
\[
      n \approx 35 \times 90 = 3150
\]
daily maxima before accounting for missing values. This is the total pooled over all winters (or all summers), not the count within a single year, since the empirical estimators pool seasonal daily maxima across the whole window. The simulation therefore uses this order of sample size as its baseline setting, while smaller values may be used in preliminary runs to assess computation time and to study how estimator performance degrades as the effective tail sample shrinks.

\subsection{Brown--Resnick benchmark process}\label{sec:simulation-br}

The main benchmark data-generating process is the Brown--Resnick max-stable process. This process is well suited to the present simulation because its extremal dependence is controlled by a variogram, which gives a direct way to vary the spatial range of extremal dependence. For two locations separated by distance $h$, the semivariogram is specified as
\[
      \gamma(h) = \left(\frac{h}{\rho}\right)^\alpha,
\]
where $\rho>0$ is a range parameter and $0<\alpha\leq 2$ controls the smoothness of the dependence decay. Larger values of $\rho$ imply slower growth of the variogram with distance and hence stronger dependence at a given distance.

For the Brown--Resnick process, the pairwise extremal coefficient is
\[
      \theta(h)
      =
      2\Phi\left(\frac{\sqrt{\gamma(h)}}{2}\right),
\]
where $\Phi$ denotes the standard normal distribution function. This expression gives the true dependence-distance curve against which the estimators can be evaluated. Since $\theta(h)\in[1,2]$, values close to one indicate strong extremal dependence, while values close to two indicate near extremal independence.

The simulation considers two regimes, labelled winter and summer. These labels are used only to mirror the empirical comparison. In the null design, both regimes have the same variogram parameters:
\[
      \rho_{\mathrm{winter}} = \rho_{\mathrm{summer}}.
\]
Therefore, there is no true seasonal difference in spatial extremal dependence. Under the alternative design, winter is assigned a larger range parameter:
\[
      \rho_{\mathrm{winter}} > \rho_{\mathrm{summer}}.
\]
This implies
\[
      \theta_{\mathrm{winter}}(h)
      <
      \theta_{\mathrm{summer}}(h)
\]
for relevant distances $h>0$, meaning that winter extremes are more spatially dependent and that this dependence persists over longer distances.

A baseline configuration uses $\alpha=1$. Under the null, both regimes may use $\rho=120\,\mathrm{km}$. Under the alternative, a representative setting is
\[
      \rho_{\mathrm{winter}} = 180\,\mathrm{km},
      \qquad
      \rho_{\mathrm{summer}} = 60\,\mathrm{km}.
\]
These values are not treated as empirical estimates. They are chosen to create a controlled and interpretable difference in dependence range.

\subsection{Estimators evaluated}\label{sec:simulation-estimators}

The simulation evaluates the same pairwise estimators used in the empirical analysis. For each Monte Carlo replication and each regime, the simulated observations are first transformed to empirical uniform margins separately by station. This matches the empirical workflow, where station-specific marginal distributions are removed before estimating spatial dependence.

Let $U_t(s_i)$ denote the transformed observation at station $s_i$ on simulated day $t$. For each pair $(i,j)$, the empirical $F$-madogram is estimated by
\[
      \widehat{\nu}_F(s_i,s_j)
      =
      \frac{1}{2n}
      \sum_{t=1}^{n}
      \left|
      U_t(s_i)-U_t(s_j)
      \right|.
\]
The corresponding madogram-based extremal coefficient estimator is
\[
      \widehat{\theta}(s_i,s_j)
      =
      \frac{1+2\widehat{\nu}_F(s_i,s_j)}
      {1-2\widehat{\nu}_F(s_i,s_j)}.
\]
The finite-level tail-dependence coefficient is estimated at thresholds $u\in\mathcal{U}$, for example $\mathcal{U}=\{0.90,0.95,0.98\}$. For a given threshold $u$, the conditional co-exceedance probability is
\[
      \widehat{\chi}_u(s_i,s_j)
      =
      \frac{
            \sum_{t=1}^{n}
            \mathbf{1}\{U_t(s_i)>u,\ U_t(s_j)>u\}
      }{
            \sum_{t=1}^{n}
            \mathbf{1}\{U_t(s_i)>u\}
      }.
\]
Because this quantity is directional in finite samples, the implementation may use a symmetrised version,
\[
      \widehat{\chi}^{\mathrm{sym}}_u(s_i,s_j)
      =
      \frac{1}{2}
      \left[
            \widehat{\Pr}\{U(s_j)>u\mid U(s_i)>u\}
            +
            \widehat{\Pr}\{U(s_i)>u\mid U(s_j)>u\}
            \right].
\]
Larger values of $\widehat{\chi}_u$ indicate stronger finite-level co-exceedance, while smaller values of $\widehat{\theta}$ indicate stronger extremal dependence.

\subsection{Monte Carlo design}\label{sec:simulation-monte-carlo}

Each Monte Carlo replication proceeds as follows. First, one winter sample and one summer sample are simulated on the fixed station set as independent draws of the max-stable field. This is an idealisation relative to the empirical data, where consecutive daily maxima within a multi-day storm are serially dependent; the simulation deliberately isolates the spatial estimation problem by treating the daily fields as independent replications, so that estimator bias and variability can be studied against a known truth without the additional complication of temporal dependence. The consequences of that residual temporal dependence are handled separately, in the empirical analysis, by the year or seasonal-block bootstrap of Section~\ref{sec:meth-uncertainty}; an extension that simulates temporally dependent daily maxima is noted in Section~\ref{sec:simulation-interpretation}. Second, each station margin is transformed to the empirical uniform scale within each regime. Third, pairwise estimates of $\theta$ and $\chi_u$ are computed for all station pairs. Fourth, the estimated winter--summer differences are computed:
\[
      \Delta_{\theta,ij}
      =
      \widehat{\theta}_{\mathrm{winter}}(s_i,s_j)
      -
      \widehat{\theta}_{\mathrm{summer}}(s_i,s_j),
\]
and
\[
      \Delta_{\chi_u,ij}
      =
      \widehat{\chi}_{u,\mathrm{winter}}(s_i,s_j)
      -
      \widehat{\chi}_{u,\mathrm{summer}}(s_i,s_j).
\]
Under the alternative design, the expected signs are
\[
      \Delta_{\theta,ij}<0
      \qquad\text{and}\qquad
      \Delta_{\chi_u,ij}>0,
\]
because winter is designed to have stronger extremal dependence. Under the null design, both differences should be centred around zero.

The pairwise estimates are also summarised over distance bins. Let $(B_1,\ldots,B_K)$ denote distance intervals, such as $0$--$50\,\mathrm{km}$, $50$--$100\,\mathrm{km}$, $100$--$150\,\mathrm{km}$, $150$--$200\,\mathrm{km}$, and above $200\,\mathrm{km}$. For each bin, the simulation reports the average estimated dependence, the average theoretical dependence, and the winter--summer difference. This converts the pairwise output into dependence-distance curves, matching the empirical analysis.

\subsection{Performance measures}\label{sec:simulation-performance}

The simulation evaluates estimator performance along several dimensions. First, for the extremal coefficient, the bias for a station pair is computed as
\[
      \operatorname{Bias}_{ij}
      =
      \mathbb{E}_{\mathrm{MC}}\!\left[
            \widehat{\theta}(s_i,s_j)
            \right]
      -
      \theta(h_{ij}),
\]
where $\mathbb{E}_{\mathrm{MC}}$ denotes the Monte Carlo average. The corresponding root mean squared error is
\[
      \operatorname{RMSE}_{ij}
      =
      \left\{
      \mathbb{E}_{\mathrm{MC}}
      \left[
            \left(
            \widehat{\theta}(s_i,s_j)-\theta(h_{ij})
            \right)^2
            \right]
      \right\}^{1/2}.
\]
These quantities show whether the madogram-based extremal coefficient estimator accurately recovers the true Brown--Resnick dependence curve.

Second, the simulation records whether the estimators recover the correct seasonal ordering. Under the alternative, the winter extremal coefficient should be lower than the summer extremal coefficient. Therefore, one diagnostic is the proportion of station pairs for which
\[
      \widehat{\theta}_{\mathrm{winter}}(s_i,s_j)
      <
      \widehat{\theta}_{\mathrm{summer}}(s_i,s_j).
\]
Analogously, for $\chi_u$, the correct ordering under the alternative is
\[
      \widehat{\chi}_{u,\mathrm{winter}}(s_i,s_j)
      >
      \widehat{\chi}_{u,\mathrm{summer}}(s_i,s_j).
\]
These diagnostics directly answer the methodological question of whether the empirical estimators are able to detect stronger winter dependence when it is present by construction.

Third, the null design is used to assess false seasonal differences. Since the two regimes have the same dependence structure under the null, the average estimated winter--summer difference should be close to zero. A systematic nonzero difference under the null would indicate that the procedure can create artificial seasonal differences even when none exist.

Fourth, the simulation checks the calibration of the uncertainty bands used in the empirical analysis. For a subset of replications, the year or seasonal-block bootstrap of Section~\ref{sec:meth-uncertainty} is applied to the simulated sample, and the empirical coverage of the nominal bootstrap interval is recorded, that is the proportion of replications in which the interval for $\theta(h_{ij})$ or for the winter--summer difference contains the known true value. Coverage close to the nominal level indicates that the bootstrap bands reported in the empirical section are trustworthy at the available sample size and station geometry; coverage well below nominal would warn that the empirical bands are too narrow. Because the simulated daily fields are independent, this exercise checks the bootstrap under the favourable case of no residual temporal dependence; the temporally dependent case is left as the extension noted in Section~\ref{sec:simulation-interpretation}.

\subsection{Interpretation for the empirical analysis}\label{sec:simulation-interpretation}

The simulation study is used to guide the empirical interpretation in two ways. First, it indicates whether the proposed estimators are reliable at the sample size and station geometry available in the KNMI data. If the estimators recover the true dependence-distance curve with small bias and acceptable variability, then the empirical winter--summer comparison can be interpreted with greater confidence. Second, the simulation clarifies the direction and magnitude of finite-sample uncertainty. In particular, it shows whether estimated seasonal differences are likely to be detectable for realistic sample sizes, or whether only large differences in dependence range can be recovered reliably.

The simulation therefore provides a bridge between the theoretical estimands and the empirical application. The empirical analysis observes only one finite sample of Dutch daily maximum wind gusts. The simulation asks whether, under comparable finite-sample conditions, the proposed methods can recover a known dependence structure. This validation step is essential for interpreting the empirical dependence-distance curves as evidence about seasonal differences in spatial extremal dependence.

A natural extension, left for future work, is to relax the independence of the simulated daily fields by introducing short-range temporal dependence between consecutive days, for example through a max-autoregressive or storm-clustered construction. Such a design would let the simulation test the year or seasonal-block bootstrap of Section~\ref{sec:meth-uncertainty} under the same multi-day storm dependence present in the KNMI data, rather than only under independent replications. This is not required to answer the main research question, which concerns whether the pairwise estimators can recover a known spatial dependence difference, but it would further strengthen the link between the simulated and empirical uncertainty quantification.